\section{Further topological restrictions}\label{app:topology-refinements}
These refinements retain information about abstract tangent splittings
and individual saturated channels. Joint saturation is governed by the
stronger covering theorem in Section~\ref{sec:general-topology}.

\subsection{What the abstract tangent splitting alone implies}
For $C^2$ factors with $P=D=M$ and $\zeta$ the identity, suppose the full mixed form
$H$ is positive definite. Each individual form
$H_i=-DA_i^*DB_i$ is positive semidefinite: the Hessian of the
nonnegative function $\ip{A_i(x)}{B_i(z)}$ at $(x,x)$ is positive
semidefinite and annihilates the tangent diagonal. Its diagonal blocks
are therefore $H_i$, and its off-diagonal blocks are $-H_i$.
If $\sum_i\rank H_i=n$ everywhere, lower semicontinuity of each rank
and constancy of their sum make the individual ranks locally constant.
Whitening by $H$ gives positive semidefinite
matrices of total rank $n$ summing to the identity. Factoring them as
$Z_iZ_i^*$ and concatenating the $n$ columns shows that each whitened
matrix is an orthogonal projection. Thus
\begin{equation}\label{eq:abstract-tangent-splitting}
 TM=\bigoplus_i E_i,\qquad E_i=\operatorname{im}(H^{-1}H_i).
\end{equation}
For exact globally smooth cone factors at capacity equality,
Theorem~\ref{thm:general-joint-cover} is stronger than the following
necessary conditions. They remain useful when only the bundle splitting
is retained.

\begin{proposition}[Ranks in a sphere tangent splitting]
\label{prop:abstract-splitting-ranks}
Suppose $T\Sph^n=\bigoplus_iE_i$, with positive ranks $r_i$, and put
$s=\rho(n+1)-1$, where
$\rho((2a+1)2^{4b+c})=8b+2^c$ for $0\leq c\leq3$.
Every subset sum $r_I=\sum_{i\in I}r_i$ satisfies
\[
 0<r_I\leq n/2\quad\Longrightarrow\quad r_I\leq s.
\]
If $s<n/3$, there is a unique rank greater than $n/2$, and it is at
least $n-s$. If all ranks are at most $c\leq n/2$, the splitting is
impossible when $s=0$; otherwise it needs at least
$\lceil n/\min\{c,s\}\rceil$ summands. For $n\geq3$, if $a$ summands
have rank one and $b$ have rank two, then $a+2b\leq s$.
\end{proposition}
\begin{proof}
For completeness, the plane-field implication can be seen from
clutching. At the endpoint $2r=n$, a splitting would make
$e(T\Sph^n)=e(E)e(F)=0$, contradicting its Euler number two.
For $0<2r<n$, both summands in a decomposition
$T\Sph^n=E^r\oplus F^{n-r}$ are orientable. Their clutching maps take
values in $SO(r)$ and $SO(n-r)$. If $r\leq n-r$, the two inclusions
in $SO(n)$ can be conjugated into the same copy of $SO(n-r)$;
the conjugations are homotopic to the identity in $SO(n)$.
The sum of the resulting homotopy classes lies in the image of
$\pi_{n-1}SO(n-r)$. Hence $T\Sph^n$ admits a reduction to $SO(n-r)$,
which supplies $r$ independent tangent vector fields. This is the
classical plane-field argument of Steenrod
\cite{Steenrod1951}. Adams's theorem bounds their number
by $s$ \cite{Adams1962}. Apply it to $\bigoplus_{i\in I}E_i$.
For $n=1$ the subset assertion is empty.

If every rank were at most $n/2$, take the first partial sum exceeding
$n/2$. The preceding sum and the last added rank are each at most $s$.
The complementary sum, when nonzero, is also at most $s$. This gives
$n\leq3s$ (and $n\leq2s$ for empty complement), a contradiction.
The complement of the unique large rank then has rank at most $s$.
The capped-count assertion follows by applying the subset bound to each
summand. For $n\geq3$, rank-one bundles are trivial and oriented
rank-two bundles are classified by their Euler class in
$H^2(\Sph^n;\mathbb Z)=0$. These summands jointly supply $a+2b$ fields.
\end{proof}

The largest-rank bound is sharp for abstract splittings: a maximal
$s$-frame and its orthogonal complement have ranks $s$ and $n-s$.
The inequality $s<n/3$ holds for $n+1\geq3$ except $n+1=4,8,16$,
as follows directly from the displayed formula for $\rho$.
For even $n$, a proper splitting is impossible already from
$e(E\oplus F)=e(E)e(F)$: both intermediate-degree Euler classes vanish,
whereas $\langle e(T\Sph^n),[\Sph^n]\rangle=2$.
For higher ranks the precise classical test is the unstable clutching
condition
\[
 \operatorname{blockdiag}_*(\alpha_1,\ldots,\alpha_k)=\tau_n
 \quad\hbox{in }\pi_{n-1}SO(n),
 \qquad\alpha_i\in\pi_{n-1}SO(r_i),
\]
where $\tau_n$ classifies $T\Sph^n$. Stable triviality of the tangent
bundle alone does not imply this condition.

\subsection{Restrictions on one everywhere saturated factor}
The normalization argument in the proof of
Theorem~\ref{thm:general-joint-cover} needs only that the chosen channel
have positive capacity $c_i$ and rank $c_i$ everywhere.
It then gives an individual submersion
$P\to\Sph^{c_i}$, even when the total capacity is larger than $\dim P$.
The next classical sphere consequence and its product extension concern
this weaker hypothesis.

\begin{proposition}[Sphere and equal-sphere-product submersions]
\label{prop:sphere-product-submersions}
A $C^1$ submersion $\Sph^n\to\Sph^r$ with $0<r<n$ is possible only
for $(n,r)=(3,2),(7,4),(15,8)$, and each pair occurs.
More generally, for $q\geq2$ and $1\leq r<kq$, a $C^1$ submersion
$(\Sph^q)^k\to\Sph^r$ requires
\begin{equation}\label{eq:product-submersion-dimensions}
 r=q\quad\hbox{or}\quad q=2r-1\ \hbox{with $r$ even}.
\end{equation}
Projection realizes the first possibility when $k\geq2$. Projection
followed by a Hopf map realizes the three corresponding possibilities
in the second case. No sufficiency assertion is made for the other
dimension pairs in~\eqref{eq:product-submersion-dimensions}.
\end{proposition}
\begin{proof}
A $C^1$ submersion from a compact source has a uniformly positive
surjectivity modulus. Smooth approximation preserves submersivity, so
Ehresmann's theorem gives a smooth fiber bundle for purposes of an
existence obstruction \cite{Ehresmann1952}. A target circle is excluded
by lifting to $\R$ and taking a maximum. For a sphere total space and
$r\geq2$, the fiber is connected. Browder's theorem gives its homotopy
type as $\Sph^j$, $j\in\{1,3,7\}$ \cite[Theorem~1]{Browder1962}.
The two-row mod-two Serre spectral sequence forces $r=j+1$: its only
possible nonzero differential killing the fiber class has length $j+1$,
and must land in the sole positive base degree $r$. The top surviving
class gives $n=2j+1$. The Hopf fibrations attain these pairs.

Now write $E=(\Sph^q)^k$. If $r\geq3$, the closed fiber $F$ is simply
connected by the homotopy exact sequence. It is orientable because the
base and total space are orientable. Both its rational homotopy and its
cohomology are finite-dimensional, so it is elliptic. The classical
formal-dimension identity \cite[Theorem~32.6]{FHT2001} is
\[
 \dim F=\sum_{j\geq2}w_jd_j(F),\qquad
 w_j=\begin{cases}j,&j\text{ odd},\\-(j-1),&j\text{ even},\end{cases}
 \quad d_j(X)=\dim_{\mathbb Q}\pi_j(X)\otimes\mathbb Q.
\]
Let $a_j$ denote the rational rank of $\pi_j(E)\to\pi_j(\Sph^r)$.
Exactness gives
$d_j(F)=d_j(E)+d_{j+1}(\Sph^r)-a_j-a_{j+1}$.
Since $w_j+w_{j-1}$ is two for odd $j$ and zero for even $j$, substitution
gives
\[
 \dim F=\begin{cases}
 kq-(r-2)-2a_r,&r\text{ odd},\\
 kq+2-r-2a_{2r-1},&r\text{ even}.
 \end{cases}
\]
Comparison with $\dim F=kq-r$ forces the displayed rank to equal one.
An odd sphere has rational homotopy only in its own degree; an even
$q$-sphere has it in degrees $q$ and $2q-1$. Thus odd $r$ forces $r=q$
if $q$ is odd. The other candidate, $r=2q-1$ for even $q$, contradicts
$2^k=\chi(E)=\chi(\Sph^r)\chi(F)=0$.
Even $r$ gives $r=q$ for even $q$, and $q=2r-1$ for odd $q$.

It remains to treat $r=2$, where the fiber need not be simply connected.
If $q\geq4$, a minimal Sullivan model of $E$ has no elements in degrees
two or three. Hence a model of $f$ sends both generators of
$(\Lambda(x_2,y_3),dy_3=x_2^2)$ to zero. The homotopy-fiber model is then
\[
 \mathcal M_E\otimes\Lambda(u_1)\otimes\mathbb Q[v_2],
 \qquad Du_1=Dv_2=0,
\]
by the standard fiber-model construction \cite[Section~15]{FHT2001}.
One can obtain it explicitly by adjoining $u_1,v_2$ to the sphere model
with $Du_1=x_2$ and $Dv_2=y_3-u_1x_2$, then pulling back along the zero
model map. Its cohomology is infinite-dimensional, contradicting the
compact fiber. Consequently $q=2$ or $3$, exactly the alternatives
in~\eqref{eq:product-submersion-dimensions} for $r=2$.
\end{proof}

In particular, one cannot apply the sphere-total-space theorem to a
product source to discard, for example, $(q,r)=(11,6)$. The rational
restriction above does not decide that case. At joint capacity equality,
however, Theorem~\ref{thm:general-joint-cover} already rules out every
smaller target, including genuine Hopf targets.

\subsection{Balanced real projector orbits}
The preceding sphere-fiber argument also restricts individual support
maps into real Grassmannians. It is useful when a single symmetric-cone
block is saturated but the whole collection is not.

\begin{proposition}[No sphere submersion onto a large balanced real orbit]
\label{prop:balanced-real-submersion}
Let $B_r=\operatorname{Gr}_{\lfloor r/2\rfloor}(\R^r)$.
For $r\geq4$ no sphere admits a $C^1$ submersion onto $B_r$ or any of
its connected finite covers. For $r=3$, the possible source dimensions
are exactly two and three. For $r=2$, they are exactly one; the maps
are circle coverings. The $r=1$ target is a point.
\end{proposition}
\begin{proof}
For $r\geq3$, a submersion from $\Sph^n$ must have $n\geq2$ and lifts
to the simply connected oriented Grassmannian
\[
 \widetilde B_r=SO(r)/(SO(p)\times SO(t)),\qquad
 p=\lfloor r/2\rfloor,\quad t=\lceil r/2\rceil.
\]
The lift is onto because its image is open and compact. Smoothing and
Ehresmann give a bundle with connected fiber. If the fiber has positive
dimension, Browder's theorem and the two-row Serre sequence give
\begin{equation}\label{eq:grass-truncated-necessary}
 H^*(\widetilde B_r;\mathbb F_2)
   =\mathbb F_2[e]/(e^{a+1}),\quad
 |e|=h\in\{2,4,8\},\quad
 pt=ah,\quad\chi(\widetilde B_r)=pt/h+1.
\end{equation}
Indeed multiplication by the transgression $e$ is an isomorphism in
all base degrees below the total-space dimension, and the fundamental
class fixes the last power.

For clarity, the required Euler characteristic can be computed directly.
The real Grassmannian has its Schubert cells indexed by $p$-subsets of
$\{1,\ldots,r\}$; the parity generating sum of their dimensions is
\[
 \chi(B_r)=
 \begin{cases}0,&p,t\text{ both odd},\\
 \binom{\lfloor r/2\rfloor}{\lfloor p/2\rfloor},&\text{otherwise}.
 \end{cases}
\]
The cell decomposition is classical \cite[Section~6]{MilnorStasheff1974}.
The formula follows by separating the subsets according to whether they contain
$r$, giving the signed Pascal recurrence. The oriented double cover has
twice this Euler characteristic. For $r=2m$, odd $m$ gives zero;
even $m\geq2$ gives $2\binom m{m/2}>m^2/2+1$, except that $m=2$
is checked directly as $4>3$. For $r=2m+1$, $m\geq4$ gives
$2\binom m{\lfloor m/2\rfloor}>m(m+1)/2+1$.
The last inequality starts at $m=4,5$ and propagates by the two-step
binomial recurrence; the even case starts at $m=4$ and uses the same
recurrence. At $m=3$ the Euler characteristic is six, whereas the
admissible values in~\eqref{eq:grass-truncated-necessary} are seven
and four. Thus the only numerical candidates are $r=3$ and $r=5$,
both with $h=2$.

For $r=5$, the map $(u,v)\mapsto[u+iv]$ identifies
$\widetilde B_5$ with the smooth three-dimensional complex quadric.
In coordinates $z_0z_4+z_1z_3+z_2^2=0$, the loci $z_0\ne0$ and
$z_0=0,z_1\ne0$ are respectively $\C^3$ and $\C^2$; the remainder
is $\mathbb{CP}^1=\C\sqcup\{\mathrm{pt}\}$.
The resulting closed filtration and compactly supported cohomology
sequence give one integral generator in each of degrees $0,2,4,6$.
Its hyperplane class $x$ is primitive: if $x=kz$ with $|k|\geq2$,
then $k^3$ would divide its cube integral, which equals two.
Poincar\'e duality therefore supplies a degree-four generator $y$ with
$\int xy=1$. The degree computation gives
$x^2=2y$. Thus the unique nonzero mod-two degree-two class has square
zero, contradicting~\eqref{eq:grass-truncated-necessary} with $a=3$.

If the fiber is zero-dimensional, the lifted map would identify a sphere
with $\widetilde B_r$. This is impossible for $r=4$, where the target is
$\Sph^2\times\Sph^2$, and for $r=5$ by the preceding cohomology.
For $r\geq6$, the homogeneous-space homotopy sequence gives
\[
 \pi_2(\widetilde B_r)
 =\ker[\mathbb Z/2\oplus\mathbb Z/2\longrightarrow\mathbb Z/2]
 =\mathbb Z/2,
\]
where the map is addition. This differs from $\pi_2(\Sph^{pt})=0$.
Finally $\widetilde B_3=\Sph^2$ gives the two stated dimensions by the
sphere-submersion result and the identity. Composing with its double
cover realizes the ordinary projective-plane target. The circle case
follows from the lifting-to-$\R$ obstruction and circle covering theory.
\end{proof}

The topology used here---plane fields, sphere fiberings, rational formal
dimension, and Grassmannian cohomology---is classical. The contribution
of this section is the passage from saturated cone contact channels to
normalized-base submersions and their joint covering, together with the
resulting formulation consequences. These statements concern one
specified globally $C^1$ slack factorization. They have different
quantifiers from the factorization of all positive linear maps through
sums of Lorentz cones studied by Aubrun, La Piana, and M\"uller-Hermes
\cite{ALM2026Lorentz}.
